\documentclass[10pt,a4paper]{amsart}
\usepackage{amsmath,amssymb,amsthm,graphicx}
% Retained preamble for the Stacks Project chapter morphisms.tex (Morphisms of Schemes).
% The chapter file pinned-original-source.tex carries only \input{preamble} on its first line, so no
% preamble text lives inside the chapter itself.  The Stacks Project uses one single shared preamble for
% every chapter; the snapshot stages it verbatim in sources/stacks-chow/preamble.tex.  The block below is
% copied mechanically, line for line, from that file's lines 14-18 (the Xy-pic setup) and lines 154-217
% (the native theorem declarations together with the mathematical macros).  Nothing was added, removed or
% redefined.  The wrapper class and the amsmath/amssymb/amsthm/graphicx base packages are supplied by the
% assembler, so the class/\externaldocument part of the shared preamble is not repeated here.
% For commutative diagrams we use Xy-pic
\usepackage[all]{xy}

% We use 2cell for 2-commutative diagrams.
\xyoption{2cell}
% Theorem environments.
%
\theoremstyle{plain}
\newtheorem{theorem}[subsection]{Theorem}
\newtheorem{proposition}[subsection]{Proposition}
\newtheorem{lemma}[subsection]{Lemma}

\theoremstyle{definition}
\newtheorem{definition}[subsection]{Definition}
\newtheorem{example}[subsection]{Example}
\newtheorem{exercise}[subsection]{Exercise}
\newtheorem{situation}[subsection]{Situation}

\theoremstyle{remark}
\newtheorem{remark}[subsection]{Remark}
\newtheorem{remarks}[subsection]{Remarks}

\numberwithin{equation}{subsection}

% Macros
%
\def\lim{\mathop{\mathrm{lim}}\nolimits}
\def\colim{\mathop{\mathrm{colim}}\nolimits}
\def\Spec{\mathop{\mathrm{Spec}}}
\def\Hom{\mathop{\mathrm{Hom}}\nolimits}
\def\Ext{\mathop{\mathrm{Ext}}\nolimits}
\def\SheafHom{\mathop{\mathcal{H}\!\mathit{om}}\nolimits}
\def\SheafExt{\mathop{\mathcal{E}\!\mathit{xt}}\nolimits}
\def\Sch{\mathit{Sch}}
\def\Mor{\mathop{\mathrm{Mor}}\nolimits}
\def\Ob{\mathop{\mathrm{Ob}}\nolimits}
\def\Sh{\mathop{\mathit{Sh}}\nolimits}
\def\NL{\mathop{N\!L}\nolimits}
\def\CH{\mathop{\mathrm{CH}}\nolimits}
\def\proetale{{pro\text{-}\acute{e}tale}}
\def\etale{{\acute{e}tale}}
\def\QCoh{\mathit{QCoh}}
\def\Ker{\mathop{\mathrm{Ker}}}
\def\Im{\mathop{\mathrm{Im}}}
\def\Coker{\mathop{\mathrm{Coker}}}
\def\Coim{\mathop{\mathrm{Coim}}}

% Boxtimes
%
\DeclareMathSymbol{\boxtimes}{\mathbin}{AMSa}{"02}

%
% Macros for moduli stacks/spaces
%
\def\QCohstack{\mathcal{QC}\!\mathit{oh}}
\def\Cohstack{\mathcal{C}\!\mathit{oh}}
\def\Spacesstack{\mathcal{S}\!\mathit{paces}}
\def\Quotfunctor{\mathrm{Quot}}
\def\Hilbfunctor{\mathrm{Hilb}}
\def\Curvesstack{\mathcal{C}\!\mathit{urves}}
\def\Polarizedstack{\mathcal{P}\!\mathit{olarized}}
\def\Complexesstack{\mathcal{C}\!\mathit{omplexes}}
% \Pic is the operator that assigns to X its picard group, usage \Pic(X)
% \Picardstack_{X/B} denotes the Picard stack of X over B
% \Picardfunctor_{X/B} denotes the Picard functor of X over B
\def\Pic{\mathop{\mathrm{Pic}}\nolimits}
\def\Picardstack{\mathcal{P}\!\mathit{ic}}
\def\Picardfunctor{\mathrm{Pic}}
\def\Deformationcategory{\mathcal{D}\!\mathit{ef}}
\title{The base change of a dominant morphism along an open morphism is dominant}
\author{The Stacks Project Authors}
\date{stacks-project:morphisms.tex:local-snapshot-20261012; GNU FDL 1.2 or later}
\begin{document}
\maketitle
\noindent\textit{Editorial note.} All delivered mathematical text is the author's own, copied from the pinned Stacks Project chapter file morphisms.tex; no statement or proof step was rewritten, reordered or removed.\par\medskip
\noindent\textit{Editorial note.} The counted statement carries the source label lemma-open-base-change-dominant; following the multi-claim convention for this source, source\_url is the source-level URL https://stacks.math.columbia.edu/ and the label is carried in claim\_key and here rather than in a per-claim tag URL.\par\medskip
\noindent\textit{Editorial note.} The counted spans contain no \textbackslash{}ref, no \textbackslash{}eqref and no citation, so no external reference stub and no bibliography excerpt are attached to this card; this was verified by scanning the two spans rather than by trusting the pool card alone.\par\medskip
\noindent\textit{Editorial note.} The proof is the author's complete argument: it does not end in a pointer, an omission or a 'follows easily', and every step that the contradiction needs is present in the counted proof span.\par\medskip
\noindent\textit{Editorial note.} The preamble is the chapter's shared Stacks Project preamble, retained as sources/stacks-morphisms/retained-preamble.tex: a line-for-line excerpt of sources/stacks-chow/preamble.tex lines 14-18 and 154-217. No macro was invented or redefined, and this card adds no preamble text of its own.\par\medskip
\noindent\textit{Allowed standard prerequisites (editorial).}\begin{itemize}
\item A morphism of schemes is dominant if and only if its image is dense, so a morphism fails to be dominant exactly when its image is contained in a proper closed subset. Same chapter, the notion is used in the counted statement at line 1336 and in the counted proof at lines 1342-1343. Supplies the form of the hypothesis and the shape of the contradiction argument: not dominant means the image lies in a closed subset B different from S'.
\item Fibre products of schemes and the universal property of the base change, in the form that points of a fibre product over a field-valued point are pairs of points with the same image. Used inside the counted proof at lines 1346-1350, where the base is changed to W = Spec(kappa(f(x))). Gives that X\_W and (S' minus B)\_W nonempty forces their fibre product over W to be nonempty, which is the contradiction that closes the proof.
\end{itemize}\medskip
\begin{lemma}
\label{lemma-open-base-change-dominant}
Let $f : X \to S$ be a dominant morphism of schemes.
Let $g : S' \to S$ be an open morphism of schemes, and let $f' : X' \to S'$
be the base change of $f$ by $g$. Then $f'$ is dominant.
\end{lemma}

\begin{proof}
If $f'$ is not dominant, then $f'(X')$ is contained in some closed
subset $B$ not equal to $S'$. Then $S'-B$ is open in $S'$ and nonempty.
Since $g$ is open, $g(S'-B)$ is nonempty and open in $S$. Since $f(X)$
is dense in $S$, there is a point $x\in X$ with $f(x) \in g(S'-B)$.
Let $W = \Spec(\kappa (f(x)))$,
and write $X_W$, $S_W$, $S'_W$ for the pullbacks
of these schemes from $S$ to $W$. Then $X_W$ and $(S'-B)_W$ are nonempty,
and so $X_W \times_W (S'-B)_W$ is nonempty, as $W$ is the spectrum
of a field. This contradicts that $f'(X \times_S S')$ is contained
in $B \subset S'$. So in fact $f'$ is dominant.
\end{proof}

\end{document}
